% Part: computability
% Chapter: machines-computations
% Section: configuration

\documentclass[../../../include/open-logic-section]{subfiles}

\begin{document}

\olfileid[hi]{cmp}{tur}{con}
\olsection{विन्यास और संगणनाएँ}

\begin{explain}
पहले सम मशीन के विन्यासों का क्रम निकालना याद करें। पट्टी, पठन-लेखन
शीर्ष और ट्यूरिंग मशीन प्रोग्राम से बनी काल्पनिक युक्ति वास्तव में केवल
यह सहज रूप से देखने की रीति है कि ट्यूरिंग मशीन की संगणना क्या होती है।
आकारिक रूप से, किसी दिए आगत पर ट्यूरिंग मशीन की संगणना को
\emph{विन्यासों} के अनुक्रम के रूप में परिभाषित कर सकते हैं---और प्रत्येक
विन्यास स्वयं प्रतीकों का एक अनुक्रम (जो संगणना के उस बिंदु पर पट्टी की
सामग्री के अनुरूप है), पठन-लेखन शीर्ष की स्थिति बताने वाली एक संख्या और
एक अवस्था है। इनकी सहायता से हम परिभाषित कर सकते हैं कि किसी दिए आगत पर
ट्यूरिंग मशीन~$M$ क्या संगणित करती है।
\end{explain}

\begin{defn}[विन्यास]
ट्यूरिंग मशीन $M = \tuple{Q, \Sigma, q_0,
\delta}$ का एक
\emph{विन्यास} त्रिक $\tuple{C, m, q}$ है, जहाँ
\begin{enumerate}
\item $C \in \Sigma^*$, $\Sigma$ के प्रतीकों का एक परिमित अनुक्रम है,
\item $m \in \Nat$, $< \len{C}$ कोई संख्या है, और
\item $q \in Q$
\end{enumerate}
सहज रूप से, अनुक्रम~$C$ पट्टी की सामग्री है (सबसे बाएँ खाने से लेकर
अंतिम अरिक्त अथवा पहले देखे गए खाने तक के सभी खानों के प्रतीक), $m$ उस
खाने की संख्या है जिसे पठन-लेखन शीर्ष पढ़ रहा है (सबसे बाएँ खाने की
संख्या $0$ से आरंभ करते हुए), और $q$ मशीन की वर्तमान अवस्था है।
\end{defn}

\begin{explain}
ट्यूरिंग मशीन का संभावित आगत प्रतीकों का एक अनुक्रम है, सामान्यतः ऐसा
अनुक्रम जो किसी संख्या को किसी रूप में कूटबद्ध करता है। ट्यूरिंग मशीन का
आरंभिक विन्यास वह विन्यास है जिसमें हम उस आगत पर काम करने के लिए मशीन
आरंभ करते हैं: पट्टी पर पट्टी-अंत चिह्नक के तुरंत बाद दाईं ओर के खानों
पर आगत लिखा होता है, पठन-लेखन शीर्ष आगत के सबसे बाएँ खाने को पढ़ता है
(अर्थात् बाएँ अंत के चिह्नक की दाईं ओर वाला खाना), और युक्ति निर्दिष्ट
आरंभिक अवस्था~$q_0$ में होती है।
\end{explain}

\begin{defn}[आरंभिक विन्यास]
$M$ का आगत $I \in \Sigma^*$ के लिए \emph{आरंभिक विन्यास} है:
\[
\tuple{\TMendtape \frown I, 1, q_0}.
\]
\end{defn}

\begin{explain}
प्रतीक~$\frown$, \emph{संलग्नन} के लिए है---आगत अक्षर-शृंखला बाएँ
अंत-चिह्नक के तुरंत बाद आरंभ होती है।
\end{explain}

\begin{defn}
हम कहते हैं कि विन्यास $\tuple{C, m, q}$ \emph{एक चरण में विन्यास
$\tuple{C', m', q'}$ देता है} ($M$ के अनुसार), तभी और केवल तभी जब
\begin{enumerate}
\item $m$वाँ प्रतीक, $C$ का, $\sigma$ है,
\item $M$ का निर्देश-समुच्चय $\delta(q, \sigma) =
  \tuple{q', \sigma', D}$ निर्दिष्ट करता है,
\item $m$वाँ प्रतीक, $C'$ का, $\sigma'$ है, और
\item
\begin{enumerate}
\item $D = L$ और $m' = m - 1$, यदि $m>0$; अन्यथा $m'=0$; अथवा
\item $D = R$ और $m' = m + 1$; अथवा
\item $D = N$ और $m' = m$,
\end{enumerate}
\item यदि $m' = \len{C}$, तो $\len{C'} = \len{C} + 1$ और $m'$वाँ
  प्रतीक, $C'$ का,~$\TMblank$ है। अन्यथा $\len{C'}=\len{C}$।
\item उन सभी $i$ के लिए जिनके लिए $i < \len{C}$ और $i \neq m$,
  $C'(i) = C(i)$,
\end{enumerate}
\end{defn}

\begin{defn}\ollabel{defn:run-output}
$M$ की \emph{आगत~$I$ पर क्रिया} विन्यासों का एक अनुक्रम $C_i$ है, जो
$M$ के विन्यास हैं, जहाँ $C_0$, $M$ का आगत~$I$ के लिए आरंभिक विन्यास
है और प्रत्येक $C_i$ एक चरण में $C_{i+1}$ देता है।

हम कहते हैं कि $M$, \emph{आगत $I$ पर $k$ चरणों के बाद रुकती है}, यदि
$C_k =
\tuple{C, m, q}$, $m$वाँ प्रतीक~$C$ का~$\sigma$ है और $\delta(q,
\sigma)$ अपरिभाषित है। उस स्थिति में $M$ का आगत~$I$ के लिए
\emph{निर्गत}~$O$ है, जहाँ $O$ प्रतीकों की ऐसी अक्षर-शृंखला है जिसका
अंत~$\TMblank$ से नहीं होता और
$C = \TMendtape \concat O \concat \TMblank^j$ किसी~$j \in \Nat$ के
लिए। ($\TMblank^j$, $j$
$\TMblank$ प्रतीकों का अनुक्रम है।)
\end{defn}

\begin{explain}
इस परिभाषा के अनुसार, निर्गत~$O$, जो~$M$ का है, सदैव~$\TMblank$ से भिन्न
किसी प्रतीक से होता है; अथवा यदि समय~$k$ पर पूरी पट्टी~$\TMblank$ से
भरी है (सबसे बाएँ~$\TMendtape$ को छोड़कर), तो $O$ रिक्त अक्षर-शृंखला
है।
\end{explain}

\end{document}
